% Part: sets-functions-relations
% Chapter: functions
% Section: functions-relations

\documentclass[../../../include/open-logic-section]{subfiles}


\begin{document}

\olfileid{sfr}{fun}{rel}

\olsection{تابعې د اړيکو په توګه}

\begin{explain}
هغه تابع چې د~$A$ !!{element}s د~$B$ له \psOblique{!!{element}s} سره
تړي، په ښکاره د $A$ او~$B$ ترمنځ يوه اړيکه ټاکي؛ يعنې هغه اړيکه
چې د $x$ او $y$ ترمنځ هله او يوازې هله وي چې $f(x) = y$ وي۔
په حقيقت کښې، که د بېلګې په توګه د رياضياتو د بنسټيزو توکو
په کمولو کښې دلچسپي لرو، ان تابع~$f$ له دې اړيکې سره، يعنې
د جوړو له يوه سټ سره، \emph{يوه ګڼلے} شو۔ بيا دا پوښتنه راپورته
کېږي: کومې اړيکې په دې ډول تابعې ټاکي؟
\end{explain}

\begin{defn}[د تابعې ګراف]فرض کړئ $f\colon A \to B$ يوه تابع ده۔
د~$f$ \emph{ګراف} هغه اړيکه $R_f \subseteq A \times B$ ده چې داسې
تعريف شوې ده:
\[
R_f = \Setabs{\tuple{x,y}}{f(x) = y}.
\]
\end{defn}

\begin{explain}
د غړو له مخې د برابرۍ په اصل، د تابعې ګراف په يوازيني ډول ټاکل
کېږي۔ سربېره پر دې، د سټونو د غړو له مخې د برابرۍ اصل سملاسي
د تابعو د قيمتونو له مخې د برابرۍ ضمني اصل هم توجيه کوي: که د
$f$ او~$g$ د تعريف ساحه او هدف سټ يو شان وي او په ټولو قيمتونو
کښې سره موافقې وي، نو هماغه يوه تابع ده۔

همدارنګه، که يوه اړيکه «د تابعې خاصيت ولري»، نو د يوې تابعې ګراف ده۔
\end{explain}

\begin{prop}\ollabel{prop:graph-function}
فرض کړئ $R \subseteq A \times B$ داسې ده چې:
\begin{enumerate}
\item که $Rxy$ او $Rxz$ وي، نو $y = z$ وي؛ او
\item د هر $x \in A$ دپاره داسې $y \in B$ شته چې $\tuple{x,
y} \in R$ وي۔
\end{enumerate}
نو $R$ د هغې تابعې $f\colon A \to B$ ګراف ده چې په دې ډول
تعريف شوې ده: $f(x) = y$ هله او يوازې هله چې $Rxy$ وي۔
\end{prop}

\begin{proof}
فرض کړئ داسې $y$ شته چې $Rxy$ وي۔ که يو بل $z \neq y$ هم واے
چې $Rxz$ واے، پر~$R$ لګېدلے شرط به مات شوے و۔ نو که داسې $y$
شته چې $Rxy$ وي، دا $y$ يوازينے دے؛ له دې کبله~$f$ په سمه توګه
تعريف شوې ده۔ په ښکاره، $R_f = R$۔
\end{proof}

\begin{explain}
هره تابع $f\colon A \to B$ ګراف لري؛ يعنې د همدې د تعريف ساحې
او هدف سټ ترمنځ يوه اړيکه، چې د $A \times B$ فرعي سټ دے او په $f(x) = y$
تعريف شوې ده۔ بل خوا، هره اړيکه~$R
\subseteq A \times B$ چې په \olref{prop:graph-function} کښې ورکړل
شوي خاصيتونه ولري، د يوې تابعې~$f \colon A \to B$ ګراف ده۔ د
تابعو او د هغو د ګرافونو د دې نږدې تړاو له کبله، تابع يوازې د
هغې د ګراف په توګه ګڼلے شو۔ په بله وينا، تابعې له ځينو اړيکو،
يعنې د مرتبو څوګونو له ځينو سټونو، سره يوې ګڼلے شو۔
\oliflabeldef{sfr:rel:ref:sec}{خو پام وکړئ چې د دې «يو ګڼلو» مطلب
هماغه دے چې په \olref[sfr][rel][ref]{sec} کښې و: دا د تابعو د
مابعدالطبيعي ماهيت په اړه ادعا نۀ ده، بلکې دا مشاهده ده چې
تابعې د ځينو سټونو په توګه \emph{ګڼل} اسانتيا لري۔ د دې اسانتيا
يو لامل دا دے چې اوس}{اوس} پر تابعو هم هغو ته ورته عمليات په
پام کښې نيولے شو چې پر اړيکو مو وکړل (وګورئ
\olref[sfr][rel][ops]{sec})۔ په ځانګړي ډول:
\end{explain}

\begin{defn}\ollabel{defn:funimage}
فرض کړئ $f \colon A \to B$ يوه تابع ده او $C\subseteq A$۔

تر~$C$ پورې د~$f$ \emph{محدودونه} هغه تابع
$\funrestrictionto{f}{C}\colon C \to B$ ده چې د ټولو $x \in C$
دپاره په $(\funrestrictionto{f}{C})(x) = f(x)$ تعريف شوې ده۔ په بله
وينا، $\funrestrictionto{f}{C} = \Setabs{\tuple{x, y} \in R_f}{x \in
C}$۔

په~$C$ باندې د~$f$ \emph{تطبيق} دا دے: $\funimage{f}{C} =
\Setabs{f(x)}{x \in C}$۔ دې ته د~$f$ لاندې د~$C$ \emph{انځور} هم وايو۔
\end{defn}

\begin{explain}
له دې تعريفونو څرګندېږي چې د هرې تابعې~$f$ دپاره
$\ran{f} = \funimage{f}{\dom{f}}$۔
\oliflabeldef{sfr:rel:ops:sec}{په \olref[sfr][rel][ops]{sec} کښې د
تعريفونو او له اړيکو سره د تابعو د يو ګڼلو له مخې، ورته
مفهومونه مو تمه لرله۔ دلته د تابعې محدودونه يوازې د ننوت
مختصات محدودوي؛ د اړيکې مخکښ تعريف دواړه مختصات محدودوي۔
خو دوه نور عمليات---معکوسونه او نسبي
حاصل ضربونه---لږ زيات تفصيل غواړي۔ هغه به په \olref[inv]{sec}
او \olref[cmp]{sec} کښې وړاندې کړو۔}{}
\end{explain}

\end{document}
